% Ad Familiares 13.2 --- headnote
% ad-familiares-13-02, around April 52 BC, Rome

\textit{Cicero to C. Memmius, written from Rome around
April 52 BC (the manuscripts: \textit{Scr. Romae, ut
videtur, m. Apr. aut paulo post a. 702}). The Memmius
addressed is C. Memmius, tribune in 54 and a
correspondent Cicero approached more than once on
behalf of clients and friends; \textit{Fam.} 13.1 and
13.3 are part of the same cluster. The subject of this
note is a small favour --- a request that Memmius let
C. Avianius Evander, a sculptor who lodges and works
in Memmius' \textit{sacrarium} (a private shrine or
chapel on his property), stay on past the Kalends of
July rather than having to move out in the middle of
extensive ongoing work.}

\textit{The letter is short and businesslike, the
standard register of the Book 13 letters of
recommendation. Cicero leans twice on the etiquette of
the genre: the appeal is qualified \textit{quod sine
tua molestia fiat} (``so far as it can be done without
trouble to you''), and the closing turn --- \textit{si
tua nihil aut non multum intersit, eo sis animo, quo
ego essem si quid tu me rogares} --- is the
characteristic Ciceronian formula of mutual obligation:
if the matter is no real concern of yours, treat my
request as you would have me treat one of yours.
Avianius Evander, the patron M. Aemilius (presumably
M. Aemilius Lepidus, the future triumvir), and the
problem of a sculptor's atelier full of half-finished
commissions all give a momentary glimpse of the texture
of patronage at the level just below the great political
moves of the year.}
